% ============================================================
%  frontmatter.tex — title page, certificate, abstract, TOC
% ============================================================

% ----------------------- TITLE PAGE -------------------------
\begin{titlepage}
\centering
\thispagestyle{empty}

{\color{navy}\rule{\textwidth}{2pt}}
\vspace{0.4cm}

{\sffamily\bfseries\LARGE\color{navy}
DEADZONE DETECTION BASED ON\\[4pt]
PACKET SNIFFING USING UAV\par}

\vspace{0.3cm}
{\color{navy}\rule{\textwidth}{2pt}}

\vspace{1.0cm}
{\itshape\large An Internship Report submitted in partial fulfillment\\
of the requirements for the award of the degree of\par}

\vspace{0.8cm}
{\sffamily\bfseries\Large\color{steel}
BACHELOR OF TECHNOLOGY\par}
\vspace{4pt}
{\sffamily\large\color{steel} in\par}
\vspace{4pt}
{\sffamily\bfseries\Large\color{steel}
COMPUTER SCIENCE AND ENGINEERING\par}

\vspace{1.2cm}
{\large Submitted by\par}
\vspace{6pt}
{\sffamily\bfseries\LARGE\color{navy} Raja Raghu Ram Pydi\par}

\vspace{0.8cm}
{\itshape\normalsize Internship carried out at\par}
\vspace{4pt}
{\bfseries\large Research Centre Imarat (RCI)\\
Defence Research \& Development Organisation (DRDO)\\
Hyderabad\par}

\vfill

\begin{tikzpicture}
  \node[draw=line, rounded corners=4pt, fill=mist, inner sep=12pt, text width=0.8\textwidth, align=center] {
    {\sffamily\bfseries\large\color{navy} DEPARTMENT OF COMPUTER SCIENCE AND ENGINEERING}\\[6pt]
    {\sffamily GITAM Institute of Technology}\\[2pt]
    {\sffamily GITAM (Deemed to be University)}\\[2pt]
    {\sffamily Visakhapatnam, India}
  };
\end{tikzpicture}

\end{titlepage}

% ----------------------- CERTIFICATE ------------------------
\chapter*{Certificate}
\addcontentsline{toc}{chapter}{Certificate}
\thispagestyle{plain}

\noindent This is to certify that the internship report entitled
\textbf{``Deadzone Detection Based on Packet Sniffing using UAV''}
is a bonafide record of the work carried out by
\textbf{Raja Raghu Ram Pydi} in partial fulfillment of the requirements for the
award of the degree of \emph{Bachelor of Technology in Computer Science and
Engineering}, GITAM (Deemed to be University). The work was undertaken during the
summer internship at Research Centre Imarat (RCI), DRDO, Hyderabad, under the
guidance named below.

\vspace{2.2cm}
\noindent
\begin{tabularx}{\textwidth}{@{}X X@{}}
\hrulefill & \hrulefill \\[2pt]
\textbf{Dr. Santanu Chatterjee} & \textbf{Prof. Thammi Reddy} \\
Scientist, Directorate of IT & Professor \& Head of Department \\
RCI, DRDO, Hyderabad & Dept. of CSE, GITAM University \\
\emph{(External Guide)} & \emph{(Internal Guide)} \\
\end{tabularx}

% ----------------------- ACKNOWLEDGEMENT --------------------
\chapter*{Acknowledgement}
\addcontentsline{toc}{chapter}{Acknowledgement}
\thispagestyle{plain}

\noindent I would like to express my special thanks and gratitude to my trainer
\textbf{Dr.~Santanu Chatterjee}, Scientist, Directorate of IT, RCI, DRDO,
Hyderabad, as well as \textbf{Prof.~Thammi Reddy}, Professor \& Head of the
Department, CSE Department, GITAM University, Visakhapatnam, who gave me the
golden opportunity to work on this wonderful project on the topic
\textbf{Deadzone Detection Based on Packet Sniffing using UAV}. This project
helped me carry out a great deal of research, and I came to learn about many new
things. I am sincerely thankful to them.

\vspace{0.8cm}
\noindent\hfill\begin{minipage}{0.45\textwidth}\raggedleft
\textbf{Raja Raghu Ram Pydi}\\
Dept. of CSE, GITAM University
\end{minipage}

% ----------------------- ABSTRACT ---------------------------
\chapter*{Abstract}
\addcontentsline{toc}{chapter}{Abstract}
\thispagestyle{plain}

\noindent As the nature of work changes, so too do the methods of automation.
Robotic process automation and intelligent systems can help organisations
improve the effectiveness of services faster and at lower cost than conventional
methods --- but with important limitations. This report studies a composite
system in which an \textbf{Unmanned Aerial Vehicle (UAV)} performs
\emph{complete-area coverage} of a target region, continuously \textbf{sniffs
wireless traffic}, and uses the resulting signal map to \textbf{detect wireless
dead-zones} (regions with no usable connectivity) and \textbf{rogue access
points}.

\vspace{6pt}
\noindent The coverage problem is cast as an optimum path-planning task over a
grid of equidistant sensing nodes whose spacing is derived from the receiver's
reception radius. A \textbf{Self-Organising Map (SOM)} neural network generates a
near-optimal traversal path and is extended to \emph{avoid obstacles} by
filtering nodes that fall inside obstacle boundaries. Flying this path, a
Raspberry-Pi--based airborne sniffer captures 802.11 management frames; gaps in
RSSI coverage mark dead-zones, while unexpected BSSIDs flag rogue access points.
A real-time analytics pipeline and an alert/feedback subsystem close the loop.

\vspace{6pt}
\noindent The system is built on a \code{Scapy}-based passive capture stack that
spans the 2.4\,GHz, 5\,GHz and 6\,GHz bands, accounts for WPA3 management-frame
protection and client MAC randomization, and performs rogue-AP classification
through a fusion of allow-listing, radio fingerprinting and
security-capability analysis. Flights follow Remote-ID / DGCA-compliant practice,
and all capture is carried out as \emph{passive}, privacy-preserving surveying
within an \emph{authorized security-testing} context.

\vspace{10pt}
\noindent\textbf{Keywords:} UAV, Packet Sniffing, Wireless Dead-Zone, Rogue
Access Point, Coverage Path Planning, Self-Organising Map, 802.11, RF Heatmap,
Raspberry Pi.

% ----------------------- TOC / LOF / LOT --------------------
\cleardoublepage
\pdfbookmark[0]{Contents}{toc}
{\color{navy}\hypersetup{linkcolor=navy}
\tableofcontents}
\cleardoublepage
{\hypersetup{linkcolor=navy}\listoffigures}
\cleardoublepage
{\hypersetup{linkcolor=navy}\listoftables}
\cleardoublepage
